% problem-000079 3 化学反应与结构综合分析
\section*{3. 化学反应与结构综合分析}
⾼⾎压是常⻅的⼼⾎管病，治疗⾼⾎压多选⽤钙离⼦拮抗剂。

\subsection*{3-1}
由苯甲醛、⼄酰⼄酸⼄酯、醋酸铵反应⽣成的化合物 $\mathrm { { \bf A } \left( C _ { 1 7 } H _ { 1 9 } N O _ { 4 } \right) }$ 是⼀种钙离⼦拮抗剂，反应式如下。画出A的结构简式

$
\mathrm{PhCHO} + 2 \mathrm{CH} _ {3} \mathrm{COCH} _ {2} \mathrm{CO} _ {2} \mathrm{CH} _ {3} + \mathrm{CH} _ {3} \mathrm{CO} _ {2} \mathrm{NH} _ {4} \xrightarrow {\mathrm{C} _ {2} \mathrm{H} _ {5} \mathrm{OH} , 8 0 ^ {\circ} \mathrm{C}} \quad \mathrm{A}
$

\subsection*{3-2}
上⾯的反应中，如果将醋酸铵转换成羟胺盐酸盐 $\left( \mathrm { N H _ { 2 } O H \cdot H C l } \right)$ ，三种物质反应先⽣成中间体B$\left( \mathrm { C _ { 1 7 } H _ { 1 9 } N O _ { 5 } } \right)$ ，进⼀步反应⽣成 $\mathrm { \bf C } \left( \mathrm { C _ { 1 7 } H _ { 1 7 } N O _ { 4 } } \right)$ 。画出B和C的结构简式。

$
\mathrm{PhCHO} + 2 \mathrm{CH} _ {3} \mathrm{COCH} _ {2} \mathrm{CO} _ {2} \mathrm{CH} _ {3} + \mathrm{NH} _ {2} \mathrm{OH} \cdot \mathrm{HCl} \xrightarrow {\mathrm{C} _ {2} \mathrm{H} _ {5} \mathrm{OH} , 8 0 ^ {\circ} \mathrm{C}} \mathrm{B} \longrightarrow \mathrm{C}
$

\subsection*{3-3}
⼄酰⼄酸⼄酯在甲醇钠的存在下与环氧⼄烷反应⽣成D，经溴化氢处理，加热⽣成E，E在氢化钠的作⽤下⽣成甲基环丙基甲酮。画出此转化过程中中间产物D和E的结构简式。

$
\mathrm{CH} _ {3} \mathrm{COCH} _ {2} \mathrm{CO} _ {2} \mathrm{CH} _ {3} \xrightarrow [ \bigtriangleup ]{\mathrm{CH} _ {3} \mathrm{ONa} , \mathrm{CH} _ {3} \mathrm{OH}} \left[ \begin{array}{c c c} & & \\ & & \end{array} \right] \xrightarrow {\mathrm{HBr}} \left[ \begin{array}{c c c} & & \\ & & \end{array} \right] \xrightarrow {\mathrm{NaH}} \mathrm{H} _ {3} \mathrm{C} ^ {\mathrm{O}}
$
